\PassOptionsToPackage{unicode=true}{hyperref} % options for packages loaded elsewhere
\PassOptionsToPackage{hyphens}{url}
%
\documentclass[]{article}
\usepackage{lmodern}
\usepackage{amssymb,amsmath}
\usepackage{ifxetex,ifluatex}
\usepackage{fixltx2e} % provides \textsubscript
\ifnum 0\ifxetex 1\fi\ifluatex 1\fi=0 % if pdftex
  \usepackage[T1]{fontenc}
  \usepackage[utf8]{inputenc}
  \usepackage{textcomp} % provides euro and other symbols
\else % if luatex or xelatex
  \usepackage{unicode-math}
  \defaultfontfeatures{Ligatures=TeX,Scale=MatchLowercase}
\fi
% use upquote if available, for straight quotes in verbatim environments
\IfFileExists{upquote.sty}{\usepackage{upquote}}{}
% use microtype if available
\IfFileExists{microtype.sty}{%
\usepackage[]{microtype}
\UseMicrotypeSet[protrusion]{basicmath} % disable protrusion for tt fonts
}{}
\IfFileExists{parskip.sty}{%
\usepackage{parskip}
}{% else
\setlength{\parindent}{0pt}
\setlength{\parskip}{6pt plus 2pt minus 1pt}
}
\usepackage{hyperref}
\hypersetup{
            pdftitle={Architectural Specification - RISC-V Core (v1)},
            pdfborder={0 0 0},
            breaklinks=true}
\urlstyle{same}  % don't use monospace font for urls
\usepackage[margin=25mm]{geometry}
\usepackage{longtable,booktabs}
% Fix footnotes in tables (requires footnote package)
\IfFileExists{footnote.sty}{\usepackage{footnote}\makesavenoteenv{longtable}}{}
\setlength{\emergencystretch}{3em}  % prevent overfull lines
\providecommand{\tightlist}{%
  \setlength{\itemsep}{0pt}\setlength{\parskip}{0pt}}
\setcounter{secnumdepth}{0}
% Redefines (sub)paragraphs to behave more like sections
\ifx\paragraph\undefined\else
\let\oldparagraph\paragraph
\renewcommand{\paragraph}[1]{\oldparagraph{#1}\mbox{}}
\fi
\ifx\subparagraph\undefined\else
\let\oldsubparagraph\subparagraph
\renewcommand{\subparagraph}[1]{\oldsubparagraph{#1}\mbox{}}
\fi

% set default figure placement to htbp
\makeatletter
\def\fps@figure{htbp}
\makeatother


\title{Architectural Specification - RISC-V Core (v1)}
\date{2026-09-01}

\begin{document}
\maketitle

\hypertarget{architectural-specification---risc-v-core-v1-rtl-snapshot-2026-09-01}{%
\section{Architectural Specification - RISC-V Core (v1, RTL snapshot
2026-09-01)}\label{architectural-specification---risc-v-core-v1-rtl-snapshot-2026-09-01}}

\hypertarget{1-document-purpose}{%
\subsection{1. Document purpose}\label{1-document-purpose}}

Define the initial (simple) specifications for the first version of a
performance-oriented RISC-V processor, to be used as a baseline for RTL
implementation, verification, and synthesis.

\hypertarget{2-first-version-goals}{%
\subsection{2. First version goals}\label{2-first-version-goals}}

\begin{itemize}
\tightlist
\item
  32-bit core for rapid ASIC flow bring-up.
\item
  Main focus: performance compatible with limited complexity.
\item
  Synthesizable and easily extensible core datapath.
\item
  Bring-up integration with cache, bus, and behavioral RAM; the sparse
  RAM model and some cache constructs are not ASIC synthesis targets in
  the current state.
\item
  No multicore support in v1.
\end{itemize}

\hypertarget{3-isa-profile}{%
\subsection{3. ISA profile}\label{3-isa-profile}}

\begin{itemize}
\tightlist
\item
  Base architecture: RV32I.
\item
  Target privilege mode: M-mode only. Privileged state is not yet
  implemented: no CSRs, trap entry/return, or interrupts are present.
\item
  Extensions in v1:

  \begin{itemize}
  \tightlist
  \item
    M (multiply/divide): not included in the first RTL iteration.
  \item
    C (compressed): not included.
  \item
    A/F/D: not included.
  \end{itemize}
\item
  Endianness: little-endian.
\item
  Instruction alignment: 32 bit.
\end{itemize}

\hypertarget{4-microarchitecture-simple-v1}{%
\subsection{4. Microarchitecture (simple
v1)}\label{4-microarchitecture-simple-v1}}

\begin{itemize}
\tightlist
\item
  3-stage in-order pipeline:

  \begin{itemize}
  \tightlist
  \item
    IF: instruction fetch
  \item
    ID/EX: decode + execute
  \item
    WB: write-back to register file
  \end{itemize}
\item
  No out-of-order execution.
\item
  No branch prediction in v1.
\item
  Two reduced AHB-Lite master ports in the core:

  \begin{itemize}
  \tightlist
  \item
    I-side read-only for instruction cache refill.
  \item
    D-side read/write for data cache refill and write-back.
  \end{itemize}
\item
  In the \texttt{riscv\_top} wrapper, the two ports are arbitrated
  toward a unified RAM with a single outstanding transaction and fixed
  D-side over I-side priority.
\item
  Control-flow handling implemented in the current RTL:

  \begin{itemize}
  \tightlist
  \item
    branch/jump resolved in ID/EX with \texttt{pc\_redirect}.
  \item
    IF/ID flush on taken redirect to eliminate the wrong-path
    instruction.
  \end{itemize}
\item
  Trap/exception handling not yet implemented (beyond
  \texttt{illegal\ instruction}).
\item
  Register file:

  \begin{itemize}
  \tightlist
  \item
    32 registers x 32 bit (x0 hardwired to 0)
  \item
    2 read ports, 1 write port.
  \end{itemize}
\item
  Reset vector hardcoded to \texttt{32\textquotesingle{}h0000\_0000} in
  the IF stage.
\end{itemize}

\hypertarget{5-functional-units}{%
\subsection{5. Functional units}\label{5-functional-units}}

\begin{itemize}
\tightlist
\item
  Integer ALU:

  \begin{itemize}
  \tightlist
  \item
    add/sub
  \item
    logic operations (and/or/xor)
  \item
    shift (sll/srl/sra)
  \item
    signed/unsigned comparisons (slt/sltu)
  \end{itemize}
\item
  Branch/jump implemented in the current RTL:

  \begin{itemize}
  \tightlist
  \item
    beq, bne, blt, bge, bltu, bgeu
  \item
    jal, jalr
  \item
    auipc (PC-relative writeback)
  \end{itemize}
\item
  Nominal load/store decode and datapath present in the current RTL:

  \begin{itemize}
  \tightlist
  \item
    lb, lh, lw, lbu, lhu
  \item
    sb, sh, sw
  \end{itemize}
\item
  End-to-end correctness of data accesses is subject to the D-cache
  limitations listed in 6.2 and 10.1.
\item
  Ecalls/ebreak/CSR/trap: not implemented in the current RTL.
\end{itemize}

\hypertarget{51-assembly-instructions-and-current-rtl-status}{%
\subsubsection{5.1 Assembly instructions and current RTL
status}\label{51-assembly-instructions-and-current-rtl-status}}

Actual implementation status in the current RTL:

\begin{itemize}
\tightlist
\item
  Implemented: \texttt{add}, \texttt{sub}, \texttt{sll}, \texttt{slt},
  \texttt{sltu}, \texttt{xor}, \texttt{srl}, \texttt{sra}, \texttt{or},
  \texttt{and}, \texttt{addi}, \texttt{slti}, \texttt{sltiu},
  \texttt{xori}, \texttt{ori}, \texttt{andi}, \texttt{slli},
  \texttt{srli}, \texttt{srai}, \texttt{lui}, \texttt{auipc},
  \texttt{jal}, \texttt{jalr}, \texttt{beq}, \texttt{bne}, \texttt{blt},
  \texttt{bge}, \texttt{bltu}, \texttt{bgeu}, \texttt{lb}, \texttt{lh},
  \texttt{lw}, \texttt{lbu}, \texttt{lhu}, \texttt{sb}, \texttt{sh},
  \texttt{sw}.
\item
  Not yet implemented: \texttt{ecall}, \texttt{ebreak}, CSR
  (\texttt{csrrw/csrrs/csrrc/csrrwi/csrrsi/csrrci}), \texttt{fence},
  \texttt{fence.i}, and other non-base-RV32I extensions of the project.
\end{itemize}

\hypertarget{register-register-arithmetic-and-logic-r-type}{%
\paragraph{Register-register arithmetic and logic
(R-type)}\label{register-register-arithmetic-and-logic-r-type}}

\begin{longtable}[]{@{}llllll@{}}
\toprule
Instruction & Description & opcode & funct3 & funct7 & Bitfield
positions\tabularnewline
\midrule
\endhead
\texttt{add\ rd,\ rs1,\ rs2} & Adds \texttt{rs1\ +\ rs2} and writes the
result to \texttt{rd}. & \texttt{0110011} & \texttt{000} &
\texttt{0000000} &
\texttt{funct7{[}31:25{]},\ rs2{[}24:20{]},\ rs1{[}19:15{]},\ funct3{[}14:12{]},\ rd{[}11:7{]},\ opcode{[}6:0{]}}\tabularnewline
\texttt{sub\ rd,\ rs1,\ rs2} & Subtracts \texttt{rs2} from \texttt{rs1}
and writes the result to \texttt{rd}. & \texttt{0110011} & \texttt{000}
& \texttt{0100000} &
\texttt{funct7{[}31:25{]},\ rs2{[}24:20{]},\ rs1{[}19:15{]},\ funct3{[}14:12{]},\ rd{[}11:7{]},\ opcode{[}6:0{]}}\tabularnewline
\texttt{sll\ rd,\ rs1,\ rs2} & Logical left shift of \texttt{rs1} by
\texttt{rs2{[}4:0{]}} bits. & \texttt{0110011} & \texttt{001} &
\texttt{0000000} &
\texttt{funct7{[}31:25{]},\ rs2{[}24:20{]},\ rs1{[}19:15{]},\ funct3{[}14:12{]},\ rd{[}11:7{]},\ opcode{[}6:0{]}}\tabularnewline
\texttt{slt\ rd,\ rs1,\ rs2} & Writes \texttt{1} to \texttt{rd} if
\texttt{rs1\ \textless{}\ rs2} signed, else \texttt{0}. &
\texttt{0110011} & \texttt{010} & \texttt{0000000} &
\texttt{funct7{[}31:25{]},\ rs2{[}24:20{]},\ rs1{[}19:15{]},\ funct3{[}14:12{]},\ rd{[}11:7{]},\ opcode{[}6:0{]}}\tabularnewline
\texttt{sltu\ rd,\ rs1,\ rs2} & Writes \texttt{1} to \texttt{rd} if
\texttt{rs1\ \textless{}\ rs2} unsigned, else \texttt{0}. &
\texttt{0110011} & \texttt{011} & \texttt{0000000} &
\texttt{funct7{[}31:25{]},\ rs2{[}24:20{]},\ rs1{[}19:15{]},\ funct3{[}14:12{]},\ rd{[}11:7{]},\ opcode{[}6:0{]}}\tabularnewline
\texttt{xor\ rd,\ rs1,\ rs2} & Bitwise XOR between \texttt{rs1} and
\texttt{rs2}. & \texttt{0110011} & \texttt{100} & \texttt{0000000} &
\texttt{funct7{[}31:25{]},\ rs2{[}24:20{]},\ rs1{[}19:15{]},\ funct3{[}14:12{]},\ rd{[}11:7{]},\ opcode{[}6:0{]}}\tabularnewline
\texttt{srl\ rd,\ rs1,\ rs2} & Logical right shift of \texttt{rs1} by
\texttt{rs2{[}4:0{]}} bits. & \texttt{0110011} & \texttt{101} &
\texttt{0000000} &
\texttt{funct7{[}31:25{]},\ rs2{[}24:20{]},\ rs1{[}19:15{]},\ funct3{[}14:12{]},\ rd{[}11:7{]},\ opcode{[}6:0{]}}\tabularnewline
\texttt{sra\ rd,\ rs1,\ rs2} & Arithmetic right shift of \texttt{rs1} by
\texttt{rs2{[}4:0{]}} bits. & \texttt{0110011} & \texttt{101} &
\texttt{0100000} &
\texttt{funct7{[}31:25{]},\ rs2{[}24:20{]},\ rs1{[}19:15{]},\ funct3{[}14:12{]},\ rd{[}11:7{]},\ opcode{[}6:0{]}}\tabularnewline
\texttt{or\ rd,\ rs1,\ rs2} & Bitwise OR between \texttt{rs1} and
\texttt{rs2}. & \texttt{0110011} & \texttt{110} & \texttt{0000000} &
\texttt{funct7{[}31:25{]},\ rs2{[}24:20{]},\ rs1{[}19:15{]},\ funct3{[}14:12{]},\ rd{[}11:7{]},\ opcode{[}6:0{]}}\tabularnewline
\texttt{and\ rd,\ rs1,\ rs2} & Bitwise AND between \texttt{rs1} and
\texttt{rs2}. & \texttt{0110011} & \texttt{111} & \texttt{0000000} &
\texttt{funct7{[}31:25{]},\ rs2{[}24:20{]},\ rs1{[}19:15{]},\ funct3{[}14:12{]},\ rd{[}11:7{]},\ opcode{[}6:0{]}}\tabularnewline
\bottomrule
\end{longtable}

\hypertarget{immediate-arithmetic-and-logic-i-type}{%
\paragraph{Immediate arithmetic and logic
(I-type)}\label{immediate-arithmetic-and-logic-i-type}}

\begin{longtable}[]{@{}llllll@{}}
\toprule
Instruction & Description & opcode & funct3 & funct7 / funct12 &
Bitfield positions\tabularnewline
\midrule
\endhead
\texttt{addi\ rd,\ rs1,\ imm} & Adds \texttt{rs1\ +\ imm} sign-extended.
& \texttt{0010011} & \texttt{000} & n/a &
\texttt{imm{[}31:20{]},\ rs1{[}19:15{]},\ funct3{[}14:12{]},\ rd{[}11:7{]},\ opcode{[}6:0{]}}\tabularnewline
\texttt{slti\ rd,\ rs1,\ imm} & Writes \texttt{1} to \texttt{rd} if
\texttt{rs1\ \textless{}\ imm} signed, else \texttt{0}. &
\texttt{0010011} & \texttt{010} & n/a &
\texttt{imm{[}31:20{]},\ rs1{[}19:15{]},\ funct3{[}14:12{]},\ rd{[}11:7{]},\ opcode{[}6:0{]}}\tabularnewline
\texttt{sltiu\ rd,\ rs1,\ imm} & Writes \texttt{1} to \texttt{rd} if
\texttt{rs1\ \textless{}\ imm} unsigned, else \texttt{0}. &
\texttt{0010011} & \texttt{011} & n/a &
\texttt{imm{[}31:20{]},\ rs1{[}19:15{]},\ funct3{[}14:12{]},\ rd{[}11:7{]},\ opcode{[}6:0{]}}\tabularnewline
\texttt{xori\ rd,\ rs1,\ imm} & Bitwise XOR between \texttt{rs1} and
\texttt{imm}. & \texttt{0010011} & \texttt{100} & n/a &
\texttt{imm{[}31:20{]},\ rs1{[}19:15{]},\ funct3{[}14:12{]},\ rd{[}11:7{]},\ opcode{[}6:0{]}}\tabularnewline
\texttt{ori\ rd,\ rs1,\ imm} & Bitwise OR between \texttt{rs1} and
\texttt{imm}. & \texttt{0010011} & \texttt{110} & n/a &
\texttt{imm{[}31:20{]},\ rs1{[}19:15{]},\ funct3{[}14:12{]},\ rd{[}11:7{]},\ opcode{[}6:0{]}}\tabularnewline
\texttt{andi\ rd,\ rs1,\ imm} & Bitwise AND between \texttt{rs1} and
\texttt{imm}. & \texttt{0010011} & \texttt{111} & n/a &
\texttt{imm{[}31:20{]},\ rs1{[}19:15{]},\ funct3{[}14:12{]},\ rd{[}11:7{]},\ opcode{[}6:0{]}}\tabularnewline
\texttt{slli\ rd,\ rs1,\ shamt} & Logical left shift of \texttt{rs1} by
\texttt{shamt} bits. & \texttt{0010011} & \texttt{001} &
\texttt{funct7=0000000} &
\texttt{funct7{[}31:25{]},\ shamt{[}24:20{]},\ rs1{[}19:15{]},\ funct3{[}14:12{]},\ rd{[}11:7{]},\ opcode{[}6:0{]}}\tabularnewline
\texttt{srli\ rd,\ rs1,\ shamt} & Logical right shift of \texttt{rs1} by
\texttt{shamt} bits. & \texttt{0010011} & \texttt{101} &
\texttt{funct7=0000000} &
\texttt{funct7{[}31:25{]},\ shamt{[}24:20{]},\ rs1{[}19:15{]},\ funct3{[}14:12{]},\ rd{[}11:7{]},\ opcode{[}6:0{]}}\tabularnewline
\texttt{srai\ rd,\ rs1,\ shamt} & Arithmetic right shift of \texttt{rs1}
by \texttt{shamt} bits. & \texttt{0010011} & \texttt{101} &
\texttt{funct7=0100000} &
\texttt{funct7{[}31:25{]},\ shamt{[}24:20{]},\ rs1{[}19:15{]},\ funct3{[}14:12{]},\ rd{[}11:7{]},\ opcode{[}6:0{]}}\tabularnewline
\bottomrule
\end{longtable}

\hypertarget{upper-immediate-and-pc-relative}{%
\paragraph{Upper immediate and
PC-relative}\label{upper-immediate-and-pc-relative}}

\begin{longtable}[]{@{}llllll@{}}
\toprule
Instruction & Description & opcode & funct3 & funct7 / note & Bitfield
positions\tabularnewline
\midrule
\endhead
\texttt{lui\ rd,\ imm20} & Loads \texttt{imm20} into the high bits of
\texttt{rd} (low bits zeroed). & \texttt{0110111} & n/a & U-type &
\texttt{imm20{[}31:12{]},\ rd{[}11:7{]},\ opcode{[}6:0{]}}\tabularnewline
\texttt{auipc\ rd,\ imm20} & Writes the value
\texttt{PC\ +\ (imm20\ \textless{}\textless{}\ 12)} to \texttt{rd}. &
\texttt{0010111} & n/a & U-type &
\texttt{imm20{[}31:12{]},\ rd{[}11:7{]},\ opcode{[}6:0{]}}\tabularnewline
\bottomrule
\end{longtable}

\hypertarget{jumps-and-branches}{%
\paragraph{Jumps and branches}\label{jumps-and-branches}}

\begin{longtable}[]{@{}llllll@{}}
\toprule
Instruction & Description & opcode & funct3 & funct7 / note & Bitfield
positions\tabularnewline
\midrule
\endhead
\texttt{jal\ rd,\ offset} & PC-relative jump and return address saved in
\texttt{rd}. & \texttt{1101111} & n/a & J-type &
\texttt{imm{[}20{]}={[}31{]},\ imm{[}10:1{]}={[}30:21{]},\ imm{[}11{]}={[}20{]},\ imm{[}19:12{]}={[}19:12{]},\ rd{[}11:7{]},\ opcode{[}6:0{]}}\tabularnewline
\texttt{jalr\ rd,\ rs1,\ imm} & Indirect jump to \texttt{rs1\ +\ imm}
with return address in \texttt{rd}. & \texttt{1100111} & \texttt{000} &
I-type &
\texttt{imm{[}31:20{]},\ rs1{[}19:15{]},\ funct3{[}14:12{]},\ rd{[}11:7{]},\ opcode{[}6:0{]}}\tabularnewline
\texttt{beq\ rs1,\ rs2,\ offset} & Branch if \texttt{rs1\ ==\ rs2}. &
\texttt{1100011} & \texttt{000} & B-type &
\texttt{imm{[}12{]}={[}31{]},\ imm{[}10:5{]}={[}30:25{]},\ rs2{[}24:20{]},\ rs1{[}19:15{]},\ funct3{[}14:12{]},\ imm{[}4:1{]}={[}11:8{]},\ imm{[}11{]}={[}7{]},\ opcode{[}6:0{]}}\tabularnewline
\texttt{bne\ rs1,\ rs2,\ offset} & Branch if \texttt{rs1\ !=\ rs2}. &
\texttt{1100011} & \texttt{001} & B-type &
\texttt{imm{[}12{]}={[}31{]},\ imm{[}10:5{]}={[}30:25{]},\ rs2{[}24:20{]},\ rs1{[}19:15{]},\ funct3{[}14:12{]},\ imm{[}4:1{]}={[}11:8{]},\ imm{[}11{]}={[}7{]},\ opcode{[}6:0{]}}\tabularnewline
\texttt{blt\ rs1,\ rs2,\ offset} & Branch if
\texttt{rs1\ \textless{}\ rs2} signed. & \texttt{1100011} & \texttt{100}
& B-type &
\texttt{imm{[}12{]}={[}31{]},\ imm{[}10:5{]}={[}30:25{]},\ rs2{[}24:20{]},\ rs1{[}19:15{]},\ funct3{[}14:12{]},\ imm{[}4:1{]}={[}11:8{]},\ imm{[}11{]}={[}7{]},\ opcode{[}6:0{]}}\tabularnewline
\texttt{bge\ rs1,\ rs2,\ offset} & Branch if
\texttt{rs1\ \textgreater{}=\ rs2} signed. & \texttt{1100011} &
\texttt{101} & B-type &
\texttt{imm{[}12{]}={[}31{]},\ imm{[}10:5{]}={[}30:25{]},\ rs2{[}24:20{]},\ rs1{[}19:15{]},\ funct3{[}14:12{]},\ imm{[}4:1{]}={[}11:8{]},\ imm{[}11{]}={[}7{]},\ opcode{[}6:0{]}}\tabularnewline
\texttt{bltu\ rs1,\ rs2,\ offset} & Branch if
\texttt{rs1\ \textless{}\ rs2} unsigned. & \texttt{1100011} &
\texttt{110} & B-type &
\texttt{imm{[}12{]}={[}31{]},\ imm{[}10:5{]}={[}30:25{]},\ rs2{[}24:20{]},\ rs1{[}19:15{]},\ funct3{[}14:12{]},\ imm{[}4:1{]}={[}11:8{]},\ imm{[}11{]}={[}7{]},\ opcode{[}6:0{]}}\tabularnewline
\texttt{bgeu\ rs1,\ rs2,\ offset} & Branch if
\texttt{rs1\ \textgreater{}=\ rs2} unsigned. & \texttt{1100011} &
\texttt{111} & B-type &
\texttt{imm{[}12{]}={[}31{]},\ imm{[}10:5{]}={[}30:25{]},\ rs2{[}24:20{]},\ rs1{[}19:15{]},\ funct3{[}14:12{]},\ imm{[}4:1{]}={[}11:8{]},\ imm{[}11{]}={[}7{]},\ opcode{[}6:0{]}}\tabularnewline
\bottomrule
\end{longtable}

\hypertarget{memory-accesses}{%
\paragraph{Memory accesses}\label{memory-accesses}}

\begin{longtable}[]{@{}llllll@{}}
\toprule
Instruction & Description & opcode & funct3 & funct7 / note & Bitfield
positions\tabularnewline
\midrule
\endhead
\texttt{lb\ rd,\ imm(rs1)} & Loads 8 bits from memory with
sign-extension. & \texttt{0000011} & \texttt{000} & I-type &
\texttt{imm{[}31:20{]},\ rs1{[}19:15{]},\ funct3{[}14:12{]},\ rd{[}11:7{]},\ opcode{[}6:0{]}}\tabularnewline
\texttt{lh\ rd,\ imm(rs1)} & Loads 16 bits from memory with
sign-extension. & \texttt{0000011} & \texttt{001} & I-type &
\texttt{imm{[}31:20{]},\ rs1{[}19:15{]},\ funct3{[}14:12{]},\ rd{[}11:7{]},\ opcode{[}6:0{]}}\tabularnewline
\texttt{lw\ rd,\ imm(rs1)} & Loads 32 bits from memory. &
\texttt{0000011} & \texttt{010} & I-type &
\texttt{imm{[}31:20{]},\ rs1{[}19:15{]},\ funct3{[}14:12{]},\ rd{[}11:7{]},\ opcode{[}6:0{]}}\tabularnewline
\texttt{lbu\ rd,\ imm(rs1)} & Loads 8 bits from memory with
zero-extension. & \texttt{0000011} & \texttt{100} & I-type &
\texttt{imm{[}31:20{]},\ rs1{[}19:15{]},\ funct3{[}14:12{]},\ rd{[}11:7{]},\ opcode{[}6:0{]}}\tabularnewline
\texttt{lhu\ rd,\ imm(rs1)} & Loads 16 bits from memory with
zero-extension. & \texttt{0000011} & \texttt{101} & I-type &
\texttt{imm{[}31:20{]},\ rs1{[}19:15{]},\ funct3{[}14:12{]},\ rd{[}11:7{]},\ opcode{[}6:0{]}}\tabularnewline
\texttt{sb\ rs2,\ imm(rs1)} & Stores 8 bits (\texttt{rs2{[}7:0{]}}) to
memory. & \texttt{0100011} & \texttt{000} & S-type &
\texttt{imm{[}11:5{]}={[}31:25{]},\ rs2{[}24:20{]},\ rs1{[}19:15{]},\ funct3{[}14:12{]},\ imm{[}4:0{]}={[}11:7{]},\ opcode{[}6:0{]}}\tabularnewline
\texttt{sh\ rs2,\ imm(rs1)} & Stores 16 bits (\texttt{rs2{[}15:0{]}}) to
memory. & \texttt{0100011} & \texttt{001} & S-type &
\texttt{imm{[}11:5{]}={[}31:25{]},\ rs2{[}24:20{]},\ rs1{[}19:15{]},\ funct3{[}14:12{]},\ imm{[}4:0{]}={[}11:7{]},\ opcode{[}6:0{]}}\tabularnewline
\texttt{sw\ rs2,\ imm(rs1)} & Stores 32 bits of \texttt{rs2} to memory.
& \texttt{0100011} & \texttt{010} & S-type &
\texttt{imm{[}11:5{]}={[}31:25{]},\ rs2{[}24:20{]},\ rs1{[}19:15{]},\ funct3{[}14:12{]},\ imm{[}4:0{]}={[}11:7{]},\ opcode{[}6:0{]}}\tabularnewline
\bottomrule
\end{longtable}

\hypertarget{system-csr-and-fence}{%
\paragraph{System, CSR, and fence}\label{system-csr-and-fence}}

\begin{longtable}[]{@{}llllll@{}}
\toprule
Instruction & Description & opcode & funct3 & funct12 / funct7 &
Bitfield positions\tabularnewline
\midrule
\endhead
\texttt{ecall} & Generates an environment call trap in M-mode. &
\texttt{1110011} & \texttt{000} & \texttt{000000000000} &
\texttt{funct12{[}31:20{]},\ rs1{[}19:15{]}=00000,\ funct3{[}14:12{]},\ rd{[}11:7{]}=00000,\ opcode{[}6:0{]}}\tabularnewline
\texttt{ebreak} & Generates a breakpoint trap in M-mode. &
\texttt{1110011} & \texttt{000} & \texttt{000000000001} &
\texttt{funct12{[}31:20{]},\ rs1{[}19:15{]}=00000,\ funct3{[}14:12{]},\ rd{[}11:7{]}=00000,\ opcode{[}6:0{]}}\tabularnewline
\texttt{csrrw\ rd,\ csr,\ rs1} & Writes \texttt{rs1} to CSR and returns
the old value in \texttt{rd}. & \texttt{1110011} & \texttt{001} &
CSR-type &
\texttt{csr{[}31:20{]},\ rs1{[}19:15{]},\ funct3{[}14:12{]},\ rd{[}11:7{]},\ opcode{[}6:0{]}}\tabularnewline
\texttt{csrrs\ rd,\ csr,\ rs1} & Sets CSR bits with \texttt{rs1},
returns the old value in \texttt{rd}. & \texttt{1110011} & \texttt{010}
& CSR-type &
\texttt{csr{[}31:20{]},\ rs1{[}19:15{]},\ funct3{[}14:12{]},\ rd{[}11:7{]},\ opcode{[}6:0{]}}\tabularnewline
\texttt{csrrc\ rd,\ csr,\ rs1} & Clears CSR bits with \texttt{rs1},
returns the old value in \texttt{rd}. & \texttt{1110011} & \texttt{011}
& CSR-type &
\texttt{csr{[}31:20{]},\ rs1{[}19:15{]},\ funct3{[}14:12{]},\ rd{[}11:7{]},\ opcode{[}6:0{]}}\tabularnewline
\texttt{csrrwi\ rd,\ csr,\ zimm} & Writes \texttt{zimm} to CSR and
returns the old value in \texttt{rd}. & \texttt{1110011} & \texttt{101}
& CSR-type &
\texttt{csr{[}31:20{]},\ zimm{[}19:15{]},\ funct3{[}14:12{]},\ rd{[}11:7{]},\ opcode{[}6:0{]}}\tabularnewline
\texttt{csrrsi\ rd,\ csr,\ zimm} & Sets CSR bits with \texttt{zimm},
returns the old value in \texttt{rd}. & \texttt{1110011} & \texttt{110}
& CSR-type &
\texttt{csr{[}31:20{]},\ zimm{[}19:15{]},\ funct3{[}14:12{]},\ rd{[}11:7{]},\ opcode{[}6:0{]}}\tabularnewline
\texttt{csrrci\ rd,\ csr,\ zimm} & Clears CSR bits with \texttt{zimm},
returns the old value in \texttt{rd}. & \texttt{1110011} & \texttt{111}
& CSR-type &
\texttt{csr{[}31:20{]},\ zimm{[}19:15{]},\ funct3{[}14:12{]},\ rd{[}11:7{]},\ opcode{[}6:0{]}}\tabularnewline
\texttt{fence} & Memory/ordering barrier (in v1 may be treated as a
functional no-op). & \texttt{0001111} & \texttt{000} & pred/succ in
immediate fields &
\texttt{fm{[}31:28{]},\ pred{[}27:24{]},\ succ{[}23:20{]},\ rs1{[}19:15{]}=00000,\ funct3{[}14:12{]},\ rd{[}11:7{]}=00000,\ opcode{[}6:0{]}}\tabularnewline
\texttt{fence.i} & Instruction fetch barrier (in v1 planned as not
implemented). & \texttt{0001111} & \texttt{001} & immediate=0 &
\texttt{imm{[}31:20{]}=000000000000,\ rs1{[}19:15{]}=00000,\ funct3{[}14:12{]},\ rd{[}11:7{]}=00000,\ opcode{[}6:0{]}}\tabularnewline
\bottomrule
\end{longtable}

v1 implementation notes:

\begin{itemize}
\tightlist
\item
  Every instruction not implemented in the current RTL generates
  \texttt{illegal\ instruction}.
\item
  Assembler pseudo-instructions (e.g., li, mv, nop, j, ret) are accepted
  by the toolchain but expanded into supported base RV32I instructions.
\end{itemize}

\hypertarget{52-risc-v-instruction-catalog-reference-for-isa-extensions}{%
\subsubsection{5.2 RISC-V instruction catalog (reference for ISA
extensions)}\label{52-risc-v-instruction-catalog-reference-for-isa-extensions}}

This chapter lists the most common standard instructions to use as an
architectural reference. For the v1 project we implement only the subset
declared in 5.1.

\hypertarget{rv32i-base-integer-isa}{%
\paragraph{RV32I (Base Integer ISA)}\label{rv32i-base-integer-isa}}

\begin{itemize}
\tightlist
\item
  lui, auipc
\item
  jal, jalr
\item
  beq, bne, blt, bge, bltu, bgeu
\item
  lb, lh, lw, lbu, lhu
\item
  sb, sh, sw
\item
  addi, slti, sltiu, xori, ori, andi, slli, srli, srai
\item
  add, sub, sll, slt, sltu, xor, srl, sra, or, and
\item
  fence, fence.i
\item
  ecall, ebreak
\item
  csrrw, csrrs, csrrc, csrrwi, csrrsi, csrrci
\end{itemize}

\hypertarget{rv64irv128i-not-a-v1-target-reference}{%
\paragraph{RV64I/RV128I (not a v1 target,
reference)}\label{rv64irv128i-not-a-v1-target-reference}}

\begin{itemize}
\tightlist
\item
  addiw, slliw, srliw, sraiw
\item
  addw, subw, sllw, srlw, sraw
\item
  ld, lwu, sd
\end{itemize}

\hypertarget{m-extension-integer-multiplydivide}{%
\paragraph{M Extension (Integer
Multiply/Divide)}\label{m-extension-integer-multiplydivide}}

\begin{itemize}
\tightlist
\item
  mul, mulh, mulhsu, mulhu
\item
  div, divu, rem, remu
\item
  Additional RV64: mulw, divw, divuw, remw, remuw
\end{itemize}

\hypertarget{a-extension-atomic}{%
\paragraph{A Extension (Atomic)}\label{a-extension-atomic}}

\begin{itemize}
\tightlist
\item
  lr.w, sc.w
\item
  amoswap.w, amoadd.w, amoxor.w, amoand.w, amoor.w
\item
  amomin.w, amomax.w, amominu.w, amomaxu.w
\item
  Additional RV64: lr.d, sc.d, amoswap.d, amoadd.d, amoxor.d, amoand.d,
  amoor.d, amomin.d, amomax.d, amominu.d, amomaxu.d
\end{itemize}

\hypertarget{f-extension-single-precision-floating-point}{%
\paragraph{F Extension (Single-Precision Floating
Point)}\label{f-extension-single-precision-floating-point}}

\begin{itemize}
\tightlist
\item
  flw, fsw
\item
  fmadd.s, fmsub.s, fnmsub.s, fnmadd.s
\item
  fadd.s, fsub.s, fmul.s, fdiv.s, fsqrt.s
\item
  fsgnj.s, fsgnjn.s, fsgnjx.s
\item
  fmin.s, fmax.s
\item
  fcvt.w.s, fcvt.wu.s
\item
  fmv.x.w
\item
  feq.s, flt.s, fle.s
\item
  fclass.s
\item
  fcvt.s.w, fcvt.s.wu
\item
  fmv.w.x
\item
  Additional RV64: fcvt.l.s, fcvt.lu.s, fcvt.s.l, fcvt.s.lu
\end{itemize}

\hypertarget{d-extension-double-precision-floating-point}{%
\paragraph{D Extension (Double-Precision Floating
Point)}\label{d-extension-double-precision-floating-point}}

\begin{itemize}
\tightlist
\item
  fld, fsd
\item
  fmadd.d, fmsub.d, fnmsub.d, fnmadd.d
\item
  fadd.d, fsub.d, fmul.d, fdiv.d, fsqrt.d
\item
  fsgnj.d, fsgnjn.d, fsgnjx.d
\item
  fmin.d, fmax.d
\item
  fcvt.s.d, fcvt.d.s
\item
  feq.d, flt.d, fle.d
\item
  fclass.d
\item
  fcvt.w.d, fcvt.wu.d, fcvt.d.w, fcvt.d.wu
\item
  Additional RV64: fcvt.l.d, fcvt.lu.d, fcvt.d.l, fcvt.d.lu
\item
  fmv.x.d, fmv.d.x (RV64)
\end{itemize}

\hypertarget{c-extension-compressed}{%
\paragraph{C Extension (Compressed)}\label{c-extension-compressed}}

\begin{itemize}
\tightlist
\item
  c.addi4spn
\item
  c.lw, c.sw
\item
  c.nop, c.addi
\item
  c.jal (RV32), c.addiw (RV64/128)
\item
  c.li, c.addi16sp, c.lui
\item
  c.srli, c.srai, c.andi
\item
  c.sub, c.xor, c.or, c.and
\item
  c.subw, c.addw (RV64/128)
\item
  c.j, c.beqz, c.bnez
\item
  c.slli
\item
  c.lwsp
\item
  c.jr, c.mv, c.ebreak, c.jalr, c.add
\item
  c.swsp
\item
  Additional RV64/128: c.ld, c.sd, c.ldsp, c.sdsp
\end{itemize}

\hypertarget{zicsr-csr-instructions}{%
\paragraph{Zicsr (CSR Instructions)}\label{zicsr-csr-instructions}}

\begin{itemize}
\tightlist
\item
  csrrw, csrrs, csrrc, csrrwi, csrrsi, csrrci
\end{itemize}

\hypertarget{zifencei-instruction-fetch-fence}{%
\paragraph{Zifencei (Instruction-Fetch
Fence)}\label{zifencei-instruction-fetch-fence}}

\begin{itemize}
\tightlist
\item
  fence.i
\end{itemize}

\hypertarget{privileged-instructions-privileged-isa-depend-on-the-supported-mode}{%
\paragraph{Privileged instructions (Privileged ISA, depend on the
supported
mode)}\label{privileged-instructions-privileged-isa-depend-on-the-supported-mode}}

\begin{itemize}
\tightlist
\item
  mret
\item
  sret
\item
  wfi
\item
  sfence.vma
\item
  hfence.vvma, hfence.gvma (hypervisor)
\end{itemize}

Notes:

\begin{itemize}
\tightlist
\item
  The set actually implemented by the core is the one declared in
  section 5.1.
\item
  This catalog serves as a reference for future evolutions (RV32IM,
  RV32IMC, RV32IMAFD, etc.).
\item
  In case of doubts about variants/edge cases, refer to the official
  RISC-V ISA specification.
\end{itemize}

\hypertarget{6-memory-subsystem-and-address-map-current-rtl-status}{%
\subsection{6. Memory subsystem and address map (current RTL
status)}\label{6-memory-subsystem-and-address-map-current-rtl-status}}

\hypertarget{61-instruction-cache}{%
\subsubsection{6.1 Instruction cache}\label{61-instruction-cache}}

\begin{itemize}
\tightlist
\item
  \texttt{riscv\_if\_stage} instantiates \texttt{riscv\_icache};
  \texttt{riscv\_imem} is not used in the current integrated datapath.
\item
  Direct-mapped, read-only cache, with default parameters:

  \begin{itemize}
  \tightlist
  \item
    128 lines.
  \item
    8 words of 32 bits per line (32 bytes).
  \item
    total capacity 4 KiB.
  \end{itemize}
\item
  Each line has a tag and a valid bit; all valid bits are cleared on
  reset.
\item
  Hit with combinational response; during a miss, the PC and the IF/ID
  register remain frozen until the refill completes.
\item
  Full-line refill via a single 256-bit transaction on the I-side
  master.
\item
  No software invalidation, no coherence with the D-cache, and no
  \texttt{fence.i} support.
\end{itemize}

\hypertarget{62-data-cache}{%
\subsubsection{6.2 Data cache}\label{62-data-cache}}

\begin{itemize}
\tightlist
\item
  \texttt{riscv\_wb\_stage} instantiates \texttt{riscv\_dcache\_ctrl},
  which wraps \texttt{set\_associative\_cache}.
\item
  Nominal default geometry:

  \begin{itemize}
  \tightlist
  \item
    8 sets x 4 ways.
  \item
    internal block of one 32-bit word.
  \item
    total data capacity 128 bytes.
  \item
    index \texttt{addr{[}4:2{]}}, tag \texttt{addr{[}31:5{]}}; way 3
    used as the eviction candidate.
  \end{itemize}
\item
  Stores are write-back: each way carries a dirty bit, and a dirty
  eviction generates a write to RAM.
\item
  \texttt{lb/lh/lbu/lhu} select byte/halfword and apply sign/zero
  extension in the controller.
\item
  \texttt{sb/sh} use a read-modify-write sequence; \texttt{sw} writes
  the word directly.
\item
  A load miss returns the word read directly from AHB but does not
  install it in the cache (\texttt{refill\_i} is tied to zero). Repeated
  loads can therefore keep missing.
\item
  Data accesses are multi-cycle. The FSM \texttt{ST\_IDLE},
  \texttt{ST\_LOAD}, \texttt{ST\_RMW\_READ}, \texttt{ST\_STORE},
  \texttt{ST\_DONE\_A}, \texttt{ST\_DONE\_B} generates
  \texttt{mem\_stall} to freeze IF and the ID/EX output registers.
\end{itemize}

D-cache limitations found in the current RTL:

\begin{itemize}
\tightlist
\item
  There are no valid bits in the ways. After reset the tags are zero,
  and an access with a zero tag can produce a false hit with zero data.
\item
  In \texttt{ST\_STORE}, \texttt{dc\_en} and \texttt{dc\_we} remain
  asserted for multiple cycles; the shift-chain structure of the ways
  can repeatedly insert the same entry and corrupt the associative
  state.
\item
  Stall is asserted starting from the state following acceptance in
  \texttt{ST\_IDLE}. ID/EX can advance on the first edge before the
  freeze takes effect, while the load destination and qualifier are not
  latched in the controller; load write-back is therefore not guaranteed
  end-to-end.
\item
  No alignment checks are present. A halfword at offset 1 or 3 is
  selected based only on bit \texttt{addr{[}1{]}}; a word ignores
  \texttt{addr{[}1:0{]}}.
\item
  The replacement policy declared in the module is not complete; the
  implementation uses a shift chain with way 3 as the eviction
  candidate.
\end{itemize}

\hypertarget{63-unified-bring-up-ram}{%
\subsubsection{6.3 Unified bring-up RAM}\label{63-unified-bring-up-ram}}

\begin{itemize}
\tightlist
\item
  \texttt{riscv\_top} instantiates \texttt{riscv\_ahb\_ram\_slave} as
  the shared backing store for the I-cache and D-cache.
\item
  32-bit byte address space; sparse storage indexed per 256-bit line.
\item
  Uninitialized locations read as zero.
\item
  Optional HEX file preload via \texttt{IMEM\_HEX\_FILE}, starting at
  \texttt{32\textquotesingle{}h0000\_0000}, for \texttt{IMEM\_WORDS}
  words (default 1024).
\item
  Read and write have a registered address phase; the model keeps
  \texttt{HREADY=1} and \texttt{HRESP=OKAY}.
\item
  A D-side write updates only the selected 32-bit lane within the line.
\item
  The sparse associative array, the \texttt{exists} method, and the
  preload \texttt{initial} block make this a simulation model, not a
  synthesizable SRAM for ASIC.
\end{itemize}

\hypertarget{64-alternativelegacy-memory-modules}{%
\subsubsection{6.4 Alternative/legacy memory
modules}\label{64-alternativelegacy-memory-modules}}

The following modules are present in the RTL folder but are not
instantiated by \texttt{riscv\_top}:

\begin{itemize}
\tightlist
\item
  \texttt{riscv\_imem}, \texttt{riscv\_dmem}: dedicated memories from
  the previous integration.
\item
  \texttt{riscv\_ahb\_imem\_slave}, \texttt{riscv\_ahb\_dmem\_slave}:
  bus wrappers for the dedicated memories.
\item
  \texttt{memory\_wrapper}, \texttt{ram}, \texttt{cdc\_sync\_2ff}:
  alternative dual-clock cache/RAM path with CDC handshake.
\end{itemize}

These modules do not describe the current integrated datapath. They also
have their own limitations: \texttt{riscv\_dmem} effectively requires
word-aligned addresses, \texttt{riscv\_ahb\_dmem\_slave} keeps the
internal request always active, and \texttt{ram} performs no range check
before indexing.

\hypertarget{7-bus-interface-current-rtl-status}{%
\subsection{7. Bus interface (current RTL
status)}\label{7-bus-interface-current-rtl-status}}

The \texttt{riscv} module exposes two reduced AHB-Lite master ports,
absorbed internally by \texttt{riscv\_top}:

\begin{itemize}
\tightlist
\item
  I-side: \texttt{HADDR}, \texttt{HTRANS}, \texttt{HWRITE=0},
  \texttt{HSIZE}, \texttt{HRDATA}, \texttt{HREADY}, \texttt{HRESP}.
\item
  D-side: same signals, with \texttt{HWRITE} and \texttt{HWDATA} for
  write-backs.
\item
  Line-wide data payload, 256 bits with default parameters.
\item
  Only \texttt{IDLE}/\texttt{NONSEQ} transfers, no bursts,
  \texttt{HPROT}, lock, or multiple outstanding transactions.
\item
  Fixed D-side over I-side priority arbitration in the top. Continuous
  D-side traffic can starve fetch.
\item
  \texttt{HRESP} reaches the caches but is not handled; no timeouts or
  traps for bus errors are present.
\end{itemize}

The interface is a custom subset compatible with the current internal
integration, not a verified general AHB-Lite implementation. In
particular, the RAM model samples \texttt{HWDATA} together with the
address phase; internal masters and slaves use the same convention, but
it does not represent the normal AHB-Lite address/data phase separation.

At the top level, \texttt{riscv\_top} exposes only \texttt{clk\_i},
\texttt{rst\_ni}, and \texttt{illegal\_instr\_o}; the bus ports are
internal to the wrapper.

\hypertarget{8-reset-clock-interrupt}{%
\subsection{8. Reset, clock, interrupt}\label{8-reset-clock-interrupt}}

\begin{itemize}
\tightlist
\item
  Single \texttt{clk\_i} clock in the active integration.
\item
  Asynchronous active-low reset used in pipeline, cache, and bus
  registers
  (\texttt{always\_ff\ @(posedge\ clk\_i\ or\ negedge\ rst\_ni)}).
\item
  Reset PC hardcoded to \texttt{32\textquotesingle{}h0000\_0000}; no
  reset vector parameter exists.
\item
  Interrupts not implemented.
\item
  \texttt{memory\_wrapper} contains a dual-clock path and 2-FF
  synchronizers, but it is not instantiated by the current top. The
  synchronizer is suitable for the single-bit handshake signals used in
  the wrapper, and does not guarantee atomicity for generic multi-bit
  buses.
\end{itemize}

\hypertarget{9-csr-traps-and-errors}{%
\subsection{9. CSR, traps, and errors}\label{9-csr-traps-and-errors}}

\begin{itemize}
\tightlist
\item
  CSRs are not implemented; there is no CSR file.
\item
  \texttt{ecall}, \texttt{ebreak}, CSR, \texttt{fence}, and
  \texttt{fence.i} are decoded as illegal instructions.
\item
  An unrecognized instruction asserts \texttt{illegal\_instr\_o} and
  inhibits write-back, but does not cause a redirect to a trap handler.
\item
  No traps are implemented for misaligned accesses, access faults,
  instruction-address-misaligned, or bus errors.
\item
  \texttt{jalr} forces bit 0 of the target to zero as required by RV32I;
  the resulting target's 32-bit alignment is not checked.
\end{itemize}

\hypertarget{10-hazard-management-and-pipeline-control}{%
\subsection{10. Hazard management and pipeline
control}\label{10-hazard-management-and-pipeline-control}}

\begin{itemize}
\tightlist
\item
  Selective data forwarding:

  \begin{itemize}
  \tightlist
  \item
    EX -\textgreater{} EX from the registered ID/EX result, with
    priority over the older bypass.
  \item
    WB -\textgreater{} EX from the final data written to the register
    file.
  \item
    match on \texttt{rs1}/\texttt{rs2}, destination different from
    \texttt{x0}, and valid write enable.
  \end{itemize}
\item
  Load-use interlock present: inserts a bubble when \texttt{load\_rd}
  matches \texttt{rs1} or \texttt{rs2} of the instruction in IF/ID.
\item
  The load-use comparison does not qualify the actual use of the source
  fields based on the opcode and can produce unnecessary conservative
  stalls.
\item
  There is no direct forwarding of load data from the D-cache to EX.
\item
  \texttt{mem\_stall} freezes PC/IF-ID and holds the ID/EX output
  registers during a multi-cycle D-cache access.
\item
  The D-cache acceptance window described in 6.2 remains an open issue:
  the existing load-use interlock does not resolve the initial loss of
  load metadata.
\item
  Branches and jumps are resolved in ID/EX without prediction. A taken
  redirect updates the PC and replaces the wrong-path instruction in
  IF/ID with a NOP, introducing a bubble.
\item
  An I-cache miss autonomously keeps PC and IF/ID frozen until
  \texttt{rvalid} goes high again.
\end{itemize}

\hypertarget{101-functional-compliance-status}{%
\subsubsection{10.1 Functional compliance
status}\label{101-functional-compliance-status}}

\begin{itemize}
\tightlist
\item
  ALU, immediates, branch/jump, and register file appear implemented
  consistently with the subset in section 5.1 based on static analysis.
\item
  Load/store decode and byte/halfword handling are present, but the
  D-cache path requires fixes and regression testing before memory
  accesses can be declared compliant.
\item
  Unimplemented instructions generate a flag, not an architectural trap.
\end{itemize}

\hypertarget{11-initial-performance-target-to-be-refined}{%
\subsection{11. Initial performance target (to be
refined)}\label{11-initial-performance-target-to-be-refined}}

\begin{itemize}
\tightlist
\item
  Initial frequency target for the first synthesis runs: 200 MHz.
\item
  Preliminary CPI target on simple code with cache hits: approximately
  1.2 - 1.8.
\item
  Taken branch: at least one bubble per flush, no prediction.
\item
  Load/store: variable latency and pipeline stall; \texttt{sb/sh} add
  the read-modify-write sequence.
\item
  I-cache misses and D-cache traffic share a single slave; D-side
  priority can increase fetch latency.
\item
  Targets must be recalibrated after the D-cache fix and with
  Genus/Innovus reports on a synthesizable memory.
\end{itemize}

\hypertarget{12-explicit-exclusions-from-v1}{%
\subsection{12. Explicit exclusions from
v1}\label{12-explicit-exclusions-from-v1}}

\begin{itemize}
\tightlist
\item
  MMU and virtual memory.
\item
  Privilege state and modes beyond the nominal machine-level datapath.
\item
  Trap handler, interrupts, and CSRs.
\item
  ISA extensions A/F/D/C/M.
\item
  Branch prediction, out-of-order execution, and multicore.
\item
  I-cache/D-cache coherence and self-modifying code support.
\item
  Complete debug hardware.
\item
  Full AHB-Lite bus interface and synthesizable integrated RAM.
\end{itemize}

\hypertarget{13-minimum-verification-requirements}{%
\subsection{13. Minimum verification
requirements}\label{13-minimum-verification-requirements}}

\begin{itemize}
\tightlist
\item
  Directed ISA tests for every instruction in section 5.1, including
  signed/unsigned and shift boundary values.
\item
  Tests for EX -\textgreater{} EX and WB -\textgreater{} EX forwarding,
  forwarding priority, and \texttt{x0} protection.
\item
  Load-use tests with instructions using one or both operands, and cases
  that should not stall.
\item
  Taken/not-taken branch/jump tests, forward/backward targets,
  wrong-path flush, and \texttt{jalr} bit 0.
\item
  I-cache tests for cold miss, hit, conflict miss, refill, wait state,
  and valid bit reset.
\item
  D-cache tests for cold access with zero tag, hit/miss on every way,
  dirty eviction, repeated loads, \texttt{sb/sh/sw}, sign extension, and
  read-modify-write.
\item
  Assertions/checks on maintaining \texttt{load\_valid} and
  \texttt{load\_rd} from request acceptance through write-back.
\item
  Simultaneous I-side/D-side starvation/arbitration tests and signal
  stability during wait states.
\item
  Negative tests for illegal instruction, misalignment, and
  \texttt{HRESP}, documenting current behavior until traps are
  implemented.
\item
  RTL code coverage and functional coverage for ISA classes, hazards,
  FSM states, hit/miss, eviction, and arbitration.
\end{itemize}

Note: the previous sequence based on dedicated IMEM/DMEM is not
sufficient to validate the current cache/AHB top. End-to-end load/store
compliance remains open until the issues in 6.2 are resolved.

\hypertarget{14-deliverables-related-to-this-specification}{%
\subsection{14. Deliverables related to this
specification}\label{14-deliverables-related-to-this-specification}}

\begin{itemize}
\tightlist
\item
  IF, ID/EX, WB pipeline and ALU/register file modules.
\item
  Direct-mapped I-cache with I-side master.
\item
  D-cache controller and set-associative cache with write-back.
\item
  I/D arbitration and unified simulation AHB RAM in the
  \texttt{riscv\_top} wrapper.
\item
  Legacy memory modules kept as alternative, inactive implementations.
\item
  Updated testbench and regression for cache, bus, stall, and
  load/store.
\item
  Initial clock/reset/IO constraints for synthesis.
\end{itemize}

\hypertarget{15-rtl-implementation-status-complete-snapshot}{%
\subsection{15. RTL implementation status (complete
snapshot)}\label{15-rtl-implementation-status-complete-snapshot}}

\hypertarget{151-active-datapath-and-control}{%
\subsubsection{15.1 Active datapath and
control}\label{151-active-datapath-and-control}}

\begin{itemize}
\tightlist
\item
  \texttt{riscv\_top.sv}: wrapper, I/D arbiter, and unified RAM.
\item
  \texttt{riscv.sv}: pipeline integration, forwarding, and interlock.
\item
  \texttt{riscv\_if\_stage.sv}: PC, IF/ID, flush, and I-cache interface.
\item
  \texttt{riscv\_id\_ex\_stage.sv}: RV32I decode, immediates, ALU
  control, branch, and registers toward WB.
\item
  \texttt{riscv\_wb\_stage.sv}: write-back mux and D-cache integration.
\item
  \texttt{riscv\_regfile.sv}: 32 x 32 bit, two asynchronous reads, one
  synchronous write.
\item
  \texttt{riscv\_alu\_addsub.sv}, \texttt{riscv\_alu\_shift.sv},
  \texttt{riscv\_alu\_logic.sv}, \texttt{riscv\_alu\_compare.sv}:
  dedicated combinational units.
\end{itemize}

\hypertarget{152-active-memory-subsystem}{%
\subsubsection{15.2 Active memory
subsystem}\label{152-active-memory-subsystem}}

\begin{itemize}
\tightlist
\item
  \texttt{riscv\_icache.sv}: direct-mapped I-cache and I-side refill.
\item
  \texttt{riscv\_dcache\_ctrl.sv}: data access control, width/sign, RMW,
  and D-side master.
\item
  \texttt{set\_associative\_cache.sv}: set selection, miss/write-back
  FSM, and RAM interface.
\item
  \texttt{cache\_set.sv}, \texttt{cache\_way.sv}: four-way storage in a
  shift chain with dirty bit.
\item
  \texttt{riscv\_ahb\_ram\_slave.sv}: unified sparse simulation backing
  store.
\end{itemize}

\hypertarget{153-alternative-modules-not-instantiated-by-the-top}{%
\subsubsection{15.3 Alternative modules not instantiated by the
top}\label{153-alternative-modules-not-instantiated-by-the-top}}

\begin{itemize}
\tightlist
\item
  \texttt{riscv\_imem.sv}, \texttt{riscv\_dmem.sv}.
\item
  \texttt{riscv\_ahb\_imem\_slave.sv},
  \texttt{riscv\_ahb\_dmem\_slave.sv}.
\item
  \texttt{memory\_wrapper.sv}, \texttt{ram.sv},
  \texttt{cdc\_sync\_2ff.sv}.
\end{itemize}

\hypertarget{154-parametrization-and-synthesis-constraints}{%
\subsubsection{15.4 Parametrization and synthesis
constraints}\label{154-parametrization-and-synthesis-constraints}}

\begin{itemize}
\tightlist
\item
  \texttt{LINE\_LENGHT} is defined as \texttt{2**LINE\_WORDS}; this
  relationship coincides with 256 bits only for the default
  \texttt{LINE\_WORDS=8}, but the general correct width would be
  \texttt{LINE\_WORDS*32}.
\item
  \texttt{ICACHE\_LINE\_WORDS} and \texttt{LINE\_WORDS} are separate
  parameters and must be kept consistent manually.
\item
  \texttt{set\_associative\_cache} extracts index and tag with fixed
  slices for 8 sets and 32-bit blocks; the declared parametrization is
  not fully general.
\item
  \texttt{cache\_set} uses an \texttt{assign\ \#1ns} delay, unsuitable
  for a portable synthesizable datapath.
\item
  \texttt{riscv\_ahb\_ram\_slave} is a sparse behavioral model and must
  be replaced with synthesizable RAM/ROM in the ASIC flow.
\end{itemize}

\hypertarget{16-open-issues-and-priorities-for-the-v1-freeze}{%
\subsection{16. Open issues and priorities for the v1
freeze}\label{16-open-issues-and-priorities-for-the-v1-freeze}}

\begin{enumerate}
\def\labelenumi{\arabic{enumi}.}
\tightlist
\item
  Fix the ID/EX-D-cache contract: assert stall in the acceptance cycle,
  or latch all metadata, including \texttt{load\_valid} and
  \texttt{load\_rd}.
\item
  Add valid bits to the D-cache and make way updates single pulses;
  verify replacement and dirty eviction.
\item
  Define and implement the load refill policy, including possible read
  allocation.
\item
  Define the behavior for misaligned accesses and add a trap or explicit
  block.
\item
  Fix the line parametrization (\texttt{LINE\_WORDS*32}) and unify the
  I-cache/system bus parameters.
\item
  Make the interface AHB-Lite compliant, or rename/document it as a
  custom internal protocol; add \texttt{HRESP} handling.
\item
  Replace sparse RAM and RTL delays with synthesizable implementations
  and define the ROM/SRAM/peripheral map.
\item
  Add cache/bus/load-store regression before declaring the RV32I subset
  fully compliant.
\item
  Decide on including the M extension, interrupt/CSR/trap strategy, and
  \texttt{fence.i} support.
\item
  Decide whether to remove legacy modules from the compile list or keep
  them in a separate test configuration.
\end{enumerate}

\end{document}
